% Part: counterfactuals
% Chapter: minimal-change-semantics
% Section: transitivity

\documentclass[../../../include/open-logic-section]{subfiles}

\begin{document}

\olfileid{cnt}{min}{cpo}

\olsection{نقيض عکس}

مادي او کلک شرطونه له خپلو نقيض عکسونو سره هم‌ارزښته دي۔
د واقعيت خلاف شرطونه داسې نۀ دي۔ د کراټزر دا بېلګه وګورئ:
\begin{quote}
  که ګويټه په 1832 کښې مړ شوے نۀ وے، نو اوس به (هم) مړ وے۔

  که ګويټه اوس مړ نۀ وے، نو په 1832 کښې به مړ شوے وے۔
\end{quote}
لومړۍ جمله رښتينې ده: انسانان سلګونه کلونه ژوند نۀ کوي۔
دويمه په څرګنده کاذبه ده: که ګويټه اوس مړ نۀ وے، نو لا به
ژوندے وے، نو په 1832 کښې مړ شوے نۀ شو کېدے۔

\begin{ex}\ollabel{ex:contraposition-counterex}
  د کروي سيمو معنٰی پوهنه نقيض عکس نامعتبر ثابتوي، يعنې $p
  \cif q \Entails/ \lnot q \cif \lnot p$ لرو۔ $p$ د دې په معنا
  وګڼئ چې «ګويټه په 1832 کښې مړ نۀ شو» او $q$ د دې په معنا
  چې «ګويټه اوس مړ دے»۔ دا په مدل $\mModel{M_1} = \tuple{W, O, V}$
  کښې بيانولے شو چې $W =
  \{w, w_1, w_2\}$، $O_w = \{\{w\}, \{w, w_1\}, \{w, w_1, w_2\}\}$،
  $V(p) = \{w_1, w_2\}$ او $V(q) = \{w, w_1\}$ لري۔ د نسخې
  څرګندونه (OLCNT-012): منجمده سرچينه د مرکز د کروي سيمو
  نظام د ټولې تابع د نوم په توګه ليکي؛ دلته د رسمي تعريف
  له مخې د مرکز د نظام نښه سمه شوې ده، ورکړل شوي سټونه عين
  دي۔ نو $w$ واقعي نړۍ ده چې ګويټه پکښې په 1832 کښې مړ شوے
  او لا مړ دے؛ $w_1$ هغه (نژدې) نړۍ ده چې ګويټه پکښې، د
  بېلګې په توګه، په 1833 کښې مړ شوے او لا مړ دے؛ او $w_2$
  هغه (ليرې) نړۍ ده چې ګويټه پکښې لا ژوندے دے۔
\begin{figure}
\begin{center}
\begin{tikzpicture}[scale=.6,layerwidth=1.5]\tiny
  \spheresystem{3}
  \begin{scope}
    \clip \propositionplot{0}{.8}{50}{4.5} -- (4.5,-4.5) -- (-4.5,-4.5) -- (-4.5,4.5) -- (4.5,4.5);
    \spherelayer{3}
  \end{scope}
  \propositionintersect{0}{2}{110}{5.5}
  \proposition{0}{.8}{50}{4.5}
  \path (0,0) node[world] {$w$};
  \spherepos{0}{2}{node[world] {$w_1$}}
  \spherepos{-50}{3}{node[world] {$w_2$}}
  \spherepos{28}{3}{node {$q$}}
  \spherepos{42}{3}{node {$\lnot q$}}
  \spherepos{-55}{4}{node {$p$}}
  \spherepos{-67}{4}{node {$\lnot p$}}
\end{tikzpicture}
\caption{د نقيض عکس ضدبېلګه}
\ollabel{fig:contraposition}
\end{center}
\end{figure}
يوه $p$-منونکې کروي سيمه $S =
\{w, w_1\}$ شته او $p \lif q$ پکښې په ټولو نړۍو کښې رښتينې
ده، نو $\mSat{M_1}{p
  \cif q}[w]$۔ د نسخې څرګندونه (OLCNT-013): منجمده سرچينه
په دې دواړو تحقق‌يي بيانونو کښې د بل مدل نوم کاروي؛ دلته
د همدې معرفي شوي مدل نوم ساتل شوے دے۔ خو $\lnot q$-منونکې
کروي سيمه $\{w, w_1,
w_2\}$ يوه نړۍ، يعنې~$w_2$، لري چې $q$ پکښې کاذب او $p$
رښتينے دے، نو $\mSat/{M_1}{\lnot q \lif \lnot p}[w_2]$۔
\end{ex}

\end{document}
